\section{Introduction}
Text-to-speech (TTS) can now clone a voice from a few seconds of reference audio.  Assistive-communication devices give a synthetic voice, ideally the user's own, to people who cannot speak. Screen readers and offline assistants run on phones and wearables. Conversational and dubbing agents must respond in real time, and because a cloned voice is biometric data, keeping it on the device rather than a server matters for privacy. These settings favor speech models that are small, fast, and on-device. Yet the quality now expected of zero-shot cloning comes from the opposite kind of model: a billion-parameter diffusion transformer (DiT) \citep{peebles2023dit} run at every step of a multi-step ODE solve \citep{anastassiou2024seedtts,chen2024f5tts,du2024cosyvoice2,xin2026longcat}. The architecture that makes personalized speech possible is too large for the devices that would use it.
 
Compression is the standard response, and post-training quantization (PTQ) is its most practical form. PTQ acts on an already-trained checkpoint, fitting low-bit weights and activations from a small calibration set with no retraining and no access to the training data, often the only feasible route for a released model. Yet every PTQ method that makes 4-bit inference work today was developed on large language models \citep{xiao2023smoothquant,frantar2023gptq,ashkboos2024quarot,sun2025flatquant} or on image and video diffusion \citep{li2023qdiffusion,zhao2025viditq,li2025svdquant}; speech synthesis has gone unstudied, and it does not behave the same way. Quantizing a TTS backbone to 4-bit weights and activations (W4A4) hurts intelligibility far more than timbre. Speaker similarity stays close to full precision, while the generated speech becomes harder to recognize. Because content and speaker identity are separate axes in speech but not in text or images, this asymmetry is speech-specific structure that methods carried over from those domains neither anticipate nor repair. Prior low-bit TTS avoids it by retraining the model \citep{kawamura2025bittts}, and the closest PTQ studies quantize non-speech audio diffusion \citep{khandelwal2025ptqaudio,vora2025ptq4adm}; to our knowledge, no PTQ study has taken a TTS diffusion transformer to W4A4 and reported what breaks.
 
We therefore treat TTS quantization as an open problem rather than a solved pipeline to re-apply. Three questions organize the study: why does 4-bit break speech, and what exactly is damaged; where does the damage live, so that it can be targeted; and can it be repaired without retraining, using only the levers PTQ already provides. To answer them we freeze a strong rotation-based backbone, FlatQuant \citep{sun2025flatquant} at its paper-best configuration, and put the effort into measurement. Every claim rests on ten independent random calibration sets per scale rather than a single draw, with strict per-utterance pairing and paired bootstrap confidence intervals for every comparison. The decisive experiments are preregistered. A W8A8 reference marks the frontier: at 8 bits intelligibility is statistically untouched at both scales and only a small timbre cost remains, so W4A4 is where speech-specific structure first appears.
 
The measurement reveals a degradation structure that is specific to speech, stable across calibration data, scales, and methods, yet impossible to read off the vision literature, whose two transferable-looking rules both reverse in speech. Concretely, we contribute:
\begin{itemize}
\item \textbf{The first systematic W4A4 PTQ study of a TTS DiT}, a degradation atlas built on ten random calibration sets per scale. It contains two phenomena not previously reported for any quantized generative audio model: a 4-bit-specific \emph{quantization bonus} (Mandarin CER improves, 10/10 sets) and a \emph{stable fragile-voice tail} (the worst decile of voices absorbs 5--10$\times$ the mean timbre loss, invisible in aggregate metrics). The structure reproduces at 3.5B and under a second W4A4 method.
\item \textbf{A new evaluation instrument}: a strict-CER taxonomy separating real acoustic errors from ASR transcription ambiguity, revealing that raw CER underestimates true intelligibility damage roughly threefold (1B $+0.13\!\to\!+0.40$) and locating it in repetition-structure stability.
\item \textbf{The first causal attribution of quantization damage in speech generation}: intelligibility damage lives in weight quantization, timbre damage in activation quantization, from same-model paired sweeps. The timbre structure is non-monotone (protecting self-attention activations is significantly \emph{worse} than protecting nothing) and constitutes a double sign-flip against the visual-DiT consensus, timestep and module axes at once.
\item \textbf{A deployable recipe}, \cfg{late-ffn}: full-precision activations only for late-step FFNs (17\% budget), the best SIM configuration at both scales, at zero intelligibility cost; and \textbf{the first controlled causal study of calibration data in TTS quantization}, whose preregistered verdict, language composition as a \emph{capacity epiphenomenon}, tells practitioners where repair does and does not live.
\end{itemize}
